\chapter{Resultados y análisis}
\label{ch:resultados}

% =======================================================================
% TODO GENERAL: este capítulo se llena con la campaña experimental.
% El protocolo base está en experimento.tex del anteproyecto (rutinas
% dominadas, ventanas de 10-60 s, criterio de viabilidad); aquí se reporta
% lo ejecutado y sus números. La macro \figpend{...} de macros.tex sirve
% como marcador de figuras pendientes.
% =======================================================================

% -----------------------------------------------------------------------
\section{Verificación funcional del clasificador}
\label{sec:res_verificacion}

% TODO:
% - Simulación (XSim): conteos del programa de prueba vs conteo manual
%   (ya existe el resultado del freeze test con filtro CSR: ARITH=6,
%   LOGIC=2, MEMORY=4, BRANCH=1, JUMP=1, FLOAT=0 — rehacer con el esquema
%   nuevo de categorías).
% - En hardware: lectura de los 12 CSR vía JTAG/GDB tras rutinas con
%   conteos conocidos; verificación de reseteo por escritura.
% - Costo en recursos FPGA del clasificador (LUTs/FFs, delta vs SoC base)
%   y efecto en timing — sale del reporte de Vivado.

% -----------------------------------------------------------------------
\section{Caracterización de coeficientes energéticos}
\label{sec:res_caracterizacion}

\subsection{Verificación de corriente máxima y configuración del PGA}
% TODO: pasada exploratoria con PGA ±4.096 V; decidir si la caracterización
% fina usa ±2.048 V (tabla de decisión del anteproyecto).

\subsection{Método de bucles dominados por categoría}
% TODO: por categoría: potencia promedio, repeticiones, dispersión,
% e_i = P*T/n_i. Tabla de coeficientes + barras con incertidumbre.
% Restar/documentar el baseline estático (p. ej. wfi).

\subsection{Método de regresión por histograma}
% TODO: M programas mixtos, matriz de conteos, mínimos cuadrados.
% Reportar R^2 y número de condición de la matriz (compromiso del
% anteproyecto, marco_teo §2.2.1). Tabla comparativa M1 vs M2.

% -----------------------------------------------------------------------
\section{Validación del modelo}
\label{sec:res_validacion}

% TODO: cargas mixtas NO usadas en caracterización; P_est vs P_medida,
% error relativo por carga y error relativo medio. Contraste explícito
% contra la hipótesis (< 10 %). Incluir casos favorables y desfavorables.

% -----------------------------------------------------------------------
\section{Discusión}
\label{sec:res_discusion}

% TODO: fuentes de error (agrupación intra-categoría, potencia estática,
% efectos de contexto/forwarding, resolución de la cadena de medición),
% comparación con la literatura (Fang, EfiMon, Jiménez) y límites de
% validez de los coeficientes (plataforma/frecuencia/bitstream).

%%% Local Variables:
%%% mode: latex
%%% TeX-master: "main"
%%% End:
